\documentclass[a4paper,12pt]{article}
\usepackage[T2A]{fontenc}
\usepackage[utf8]{inputenc}
\usepackage[russian]{babel}
\usepackage{amsmath,amssymb,mathtools}
\renewcommand{\familydefault}{\sfdefault}
\usepackage{geometry}
\geometry{top=24mm,bottom=22mm,left=25mm,right=25mm}
\usepackage{microtype}
\usepackage{icomma}
\usepackage{booktabs,tabularx,longtable}
\usepackage{enumitem}
\usepackage{fancyhdr}
\usepackage{parskip}
\usepackage{setspace}
\usepackage{xcolor}
\usepackage{tikz}
\usetikzlibrary{calc,arrows.meta,positioning,shapes.geometric}
\usepackage{pgfplots}
\pgfplotsset{compat=1.18}
\usepackage{hyperref}
\hypersetup{hidelinks,pdfauthor={Лёвин Артём Александрович},pdftitle={Чек-лист. Линейная функция. Кирилл Зиновьев}}
\definecolor{accent}{HTML}{007F83}
\definecolor{deepaccent}{HTML}{005A61}
\definecolor{softgray}{HTML}{575D60}
\definecolor{warmgray}{HTML}{8B9294}
\definecolor{ink}{HTML}{2A2E32}
\color{ink}
\onehalfspacing
\setlist[itemize]{leftmargin=6mm,itemsep=2pt,topsep=3pt}
\setlist[enumerate]{leftmargin=8mm,itemsep=7pt,topsep=4pt}
\setlength{\parskip}{7pt}
\setlength{\parindent}{0pt}
\renewcommand{\arraystretch}{1.24}
\pagestyle{fancy}
\fancyhf{}
\renewcommand{\headrulewidth}{0pt}
\fancyhead[L]{\footnotesize\color{softgray}Линейная функция и её график}
\fancyhead[R]{\footnotesize\color{softgray}Кирилл Зиновьев}
\fancyfoot[C]{\footnotesize\color{softgray}\thepage}
\setlength{\headheight}{15pt}
\newcommand{\topic}[1]{\refstepcounter{section}\vspace{10pt}\noindent{\large\color{deepaccent}\thesection.\enspace #1}\par\vspace{4pt}{\color{accent}\hrule height .48pt}\vspace{9pt}}
\newcommand{\smalltopic}[1]{\vspace{5pt}\noindent\textcolor{deepaccent}{\bfseries #1}\par\vspace{2pt}}
\newcommand{\ex}[1]{\noindent\textcolor{deepaccent}{\bfseries Пример.}\ #1}
\newcommand{\keyidea}[1]{\noindent\textcolor{deepaccent}{\bfseries Главное.}\ #1}
\newcommand{\warning}[1]{\noindent\textcolor{deepaccent}{\bfseries Проверьте себя.}\ #1}
\begin{document}
\thispagestyle{empty}
\begin{center}
\vspace*{22mm}
{\color{deepaccent}\large МАТЕМАТИКА}\par
\vspace{17mm}
{\LARGE Чек-лист}\par
\vspace{9mm}
{\color{accent}\LARGE Линейная функция:\par графики, коэффициенты\par и точки пересечения}\par
\vspace{17mm}
{\normalsize Построение прямой, вычисление координат\par и анализ графика}\par
\vfill
{\large Кирилл Зиновьев}\par
\vspace{11mm}
{14 ноября 2026 г.}\par
\vspace{16mm}
{\small Лёвин Артём Александрович\par эксклюзивно для Кирилла Зиновьева}\par
\vspace{15mm}
\begin{tikzpicture}[x=1cm,y=1cm]
  \draw[accent, line width=.75pt] (0,0) .. controls (1.35,.52) and (2.9,-.48) .. (4.3,0);
  \draw[accent, line width=.75pt] (4.3,0) .. controls (5.8,.52) and (7.15,-.48) .. (8.6,0);
\end{tikzpicture}
\vspace{2mm}
\end{center}
\clearpage
\setcounter{page}{1}

\topic{Что называется линейной функцией}
Линейная функция задаётся формулой
\[
\boxed{y=kx+b},
\]
где $x$ --- аргумент, $y$ --- значение функции, $k$ --- \textbf{угловой коэффициент}, $b$ --- \textbf{свободный коэффициент}. Числа $k$ и $b$ могут быть положительными, отрицательными или равными нулю.

\keyidea{График функции $y=kx+b$ --- прямая. При $k=0$ она горизонтальна: $y=b$.}

\topic{Как найти значение функции}
Если задано значение $x$, подставьте его в формулу вместо переменной $x$. Отрицательные числа заключайте в скобки.

\ex{Для функции $y=-3x+1{,}5$ найдём значение при $x=-1{,}5$:}
\[
\begin{aligned}
y(-1{,}5)&=-3\cdot(-1{,}5)+1{,}5\\
&=4{,}5+1{,}5=6.
\end{aligned}
\]
Для той же функции:
\[
 y(2{,}5)=-3\cdot2{,}5+1{,}5=-6,
 \qquad y(4)=-12+1{,}5=-10{,}5.
\]

\warning{Умножение отрицательного числа на положительное даёт отрицательный результат. В конце проверьте знак каждого слагаемого.}

\topic{Обратная задача: найти аргумент по значению функции}
Если известно $y$, поставьте указанное число в \textbf{левую часть формулы} и решите полученное уравнение относительно $x$.

\ex{Найдём $x$, если $y=-4{,}5$ и $y=-3x+1{,}5$:}
\[
-4{,}5=-3x+1{,}5
\ \Longrightarrow\ -6=-3x
\ \Longrightarrow\ \boxed{x=2}.
\]
Проверка: $-3\cdot2+1{,}5=-4{,}5$.

\smalltopic{Полезная запись}
$y(a)$ обозначает значение функции при $x=a$. Например, $y(-1{,}5)=6$. Такую запись удобно использовать при последовательной подстановке чисел.

\clearpage
\topic{Построение прямой по двум точкам}
Чтобы построить график $y=kx+b$, достаточно найти \textbf{две различные точки} этой прямой.

\smalltopic{Порядок построения}
\begin{enumerate}
\item Выберите два разных значения $x$ (обычно удобны $0$ и $1$).
\item Для каждого значения вычислите $y$ по формуле.
\item Составьте таблицу, отметьте точки $(x;y)$ на координатной плоскости.
\item Проведите через обе точки \textbf{прямую} и продолжите её в обе стороны.
\end{enumerate}

\ex{Построим $y=-2x+1$.}
\[
\begin{array}{c|cc}
x&0&1\\\hline
y=-2x+1&1&-1
\end{array}
\qquad A(0;1),\quad B(1;-1).
\]
\begin{center}
\begin{tikzpicture}
\begin{axis}[
  width=.82\linewidth,height=65mm,
  axis lines=middle,xmin=-1.3,xmax=2.6,ymin=-3.6,ymax=3.5,
  xtick={-1,0,1,2},ytick={-3,-2,-1,0,1,2,3},
  grid=both,grid style={gray!14},minor tick num=0,
  xlabel={$x$},ylabel={$y$},xlabel style={at={(axis description cs:1,.44)},anchor=west},
  ylabel style={at={(axis description cs:.3,1)},anchor=south},
  clip=false,tick label style={font=\small},
]
\addplot[domain=-1.2:2.35,samples=2,thick,color=accent] {-2*x+1};
\addplot[only marks,mark=*,mark size=2.2pt,color=deepaccent] coordinates {(0,1) (1,-1)};
\node[font=\small,above right] at (axis cs:0,1) {$A(0;1)$};
\node[font=\small,below right] at (axis cs:1,-1) {$B(1;-1)$};
\end{axis}
\end{tikzpicture}
\end{center}
\warning{Первая координата точки --- $x$ (абсцисса), вторая --- $y$ (ордината). Точки служат опорой для построения всей прямой.}

\topic{Что показывает угловой коэффициент $k$}
\begin{center}
\begin{tabularx}{\linewidth}{@{}lX@{}}
\toprule
Условие & Как ведёт себя прямая при движении \textbf{слева направо}\\
\midrule
$k>0$ & Функция возрастает: прямая идёт снизу вверх.\\
$k<0$ & Функция убывает: прямая идёт сверху вниз.\\
$k=0$ & Функция постоянна: прямая горизонтальна.\\
\bottomrule
\end{tabularx}
\end{center}
\clearpage

\topic{Что показывает свободный коэффициент $b$}
Подставьте $x=0$ в $y=kx+b$:
\[
y=k\cdot0+b=b.
\]
Значит, график пересекает ось $Oy$ в точке \(\boxed{(0;b)}\).

\begin{itemize}
\item Если $b>0$, точка пересечения с $Oy$ находится \textbf{выше} начала координат.
\item Если $b<0$, точка находится \textbf{ниже} начала координат.
\item Если $b=0$, прямая проходит через начало координат.
\end{itemize}

\smalltopic{Сравнение двух прямых}
У функций $y=2x+1$ и $y=-2x+1$ коэффициент $b$ одинаков. Поэтому обе прямые проходят через точку $(0;1)$. Положительное и отрицательное значения $k$ задают противоположные направления.
\begin{center}
\begin{tikzpicture}
\begin{axis}[
 width=.84\linewidth,height=50mm,axis lines=middle,
 xmin=-1.55,xmax=1.55,ymin=-2.3,ymax=4.3,
 xtick={-1,0,1},ytick={-2,-1,0,1,2,3,4},
 grid=both,grid style={gray!12},
 xlabel={$x$},ylabel={$y$},
 legend pos=north east,legend style={font=\small,draw=none,fill=none},
 tick label style={font=\small},
]
\addplot[domain=-1.5:1.5,samples=2,very thick,color=accent] {2*x+1};
\addlegendentry{$y=2x+1$}
\addplot[domain=-1.5:1.5,samples=2,thick,color=softgray,dashed] {-2*x+1};
\addlegendentry{$y=-2x+1$}
\addplot[only marks,mark=*,mark size=2.1pt,color=deepaccent] coordinates {(0,1)};
\end{axis}
\end{tikzpicture}
\end{center}
\keyidea{$k$ определяет наклон прямой; $b$ определяет её пересечение с осью $Oy$. По этим двум признакам можно быстро выполнить эскиз.}

\topic{Эскиз по знакам $k$ и $b$}
В каждой из четырёх ситуаций сначала определите направление прямой по $k$, затем найдите, выше или ниже нуля она пересекает ось $Oy$, по $b$.
\begin{center}
\renewcommand{\arraystretch}{1.45}
\begin{tabularx}{\linewidth}{@{}llX@{}}
\toprule
$k$&$b$&Эскиз (словесное описание)\\
\midrule
$k>0$&$b>0$&Возрастает; пересекает $Oy$ выше нуля.\\
$k>0$&$b<0$&Возрастает; пересекает $Oy$ ниже нуля.\\
$k<0$&$b>0$&Убывает; пересекает $Oy$ выше нуля.\\
$k<0$&$b<0$&Убывает; пересекает $Oy$ ниже нуля.\\
\bottomrule
\end{tabularx}
\end{center}
\clearpage

\topic{Как найти точку пересечения двух прямых}
В точке пересечения графики имеют \textbf{одинаковые значения $x$ и $y$}. Поэтому правые части формул приравнивают.

\ex{Найдём точку пересечения $y=4x+1$ и $y=-2x+4$:}
\[
4x+1=-2x+4
\ \Longrightarrow\ 6x=3
\ \Longrightarrow\ \boxed{x=0{,}5}.
\]
Теперь найдём ординату, подставив $x=0{,}5$:
\[
 y=4\cdot0{,}5+1=3.
\]
Проверим вторую формулу: $y=-2\cdot0{,}5+4=3$.

\keyidea{Точка пересечения: $\boxed{(0{,}5;3)}$. Проверка требует получить одинаковый $y$ по обеим формулам.}

\smalltopic{Второй пример: дробная абсцисса}
\[
\begin{aligned}
y_1&=3x-4,\\
y_2&=-2x+5.
\end{aligned}
\qquad
3x-4=-2x+5
\ \Longrightarrow\ x=\frac95=1{,}8.
\]
\[
 y_1=3\cdot1{,}8-4=1{,}4,
 \qquad y_2=-2\cdot1{,}8+5=1{,}4.
\]
Ответ: $\boxed{(1{,}8;1{,}4)}$.

\smalltopic{Сколько может быть точек пересечения?}
Пусть прямые заданы формулами $y=k_1x+b_1$ и $y=k_2x+b_2$:
\begin{itemize}
\item при $k_1\ne k_2$ существует \textbf{одна} точка пересечения;
\item при $k_1=k_2$ и $b_1\ne b_2$ прямые параллельны, общих точек нет;
\item при $k_1=k_2$ и $b_1=b_2$ прямые совпадают, общих точек бесконечно много.
\end{itemize}

\warning{Уравнение для пересечения даёт абсциссу. Для полного ответа найдите ординату, а затем проверьте её по второй формуле.}

\clearpage
\thispagestyle{fancy}
\begin{center}
{\large\color{deepaccent}Алгоритм: точка пересечения двух графиков}
\end{center}
\vspace{6mm}
\begin{center}
\begin{tikzpicture}[
 every node/.style={font=\small,align=center},
 beginend/.style={ellipse,draw=deepaccent,line width=.8pt,minimum width=42mm,minimum height=10mm,inner sep=4pt},
 action/.style={rectangle,rounded corners=2pt,draw=accent,line width=.75pt,text width=84mm,minimum height=13mm,inner sep=5pt},
 decision/.style={diamond,aspect=2.45,draw=deepaccent,line width=.8pt,text width=52mm,inner sep=3pt},
 edge/.style={-{Stealth[length=2.6mm]},line width=.8pt,color=deepaccent},
 node distance=11mm,
]
\node[beginend] (start) {Даны две формулы $y=f(x)$ и $y=g(x)$};
\node[action,below=of start] (equal) {Приравняйте правые части: $f(x)=g(x)$};
\node[action,below=of equal] (solve) {Решите уравнение и найдите абсциссу $x$};
\node[action,below=of solve] (value) {Подставьте $x$ в обе формулы и вычислите $y_1$, $y_2$};
\node[decision,below=13mm of value] (check) {$y_1=y_2$?};
\node[action,below=14mm of check] (answer) {Запишите точку $(x;y)$};
\node[beginend,below=of answer] (finish) {Готово};
\node[action,right=21mm of check,text width=40mm] (redo) {Проверьте\ знаки и вычисления};
\draw[edge] (start)--(equal);
\draw[edge] (equal)--(solve);
\draw[edge] (solve)--(value);
\draw[edge] (value)--(check);
\draw[edge] (check)--node[left,font=\footnotesize]{да}(answer);
\draw[edge] (answer)--(finish);
\draw[edge] (check.east)--node[above,font=\footnotesize]{нет}(redo.west);
\draw[edge] (redo.north)|- (value.east);
\end{tikzpicture}
\end{center}
\clearpage

\topic{Типичные ошибки и короткая памятка}
\begin{itemize}
\item \textbf{Пропущенные скобки.} При подстановке отрицательного $x$ обязательно пишите, например, $-3\cdot(-1{,}5)$.
\item \textbf{Ошибка знака.} $-7{,}5+1{,}5=-6$.
\item \textbf{Перепутаны координаты.} Запись $(x;y)$ начинается с абсциссы.
\item \textbf{Построена только часть прямой.} Через найденные точки проведите прямую в обе стороны.
\item \textbf{Забыли найти $y$.} После решения $f(x)=g(x)$ нужно вычислить ординату.
\item \textbf{Не выполнена проверка.} Подставьте найденный $x$ в обе формулы; результаты должны совпасть.
\end{itemize}
\smalltopic{Итоговая памятка}
\begin{center}
\begin{tabularx}{\linewidth}{@{}p{.41\linewidth}X@{}}
\toprule
Сигнал в задаче & Полезное действие\\
\midrule
«Найдите $y$, если $x=\ldots$» & Подставьте $x$ в правую часть формулы.\\
«Найдите $x$, если $y=\ldots$» & Подставьте $y$ и решите уравнение.\\
«Постройте график» & Возьмите два разных $x$, вычислите $y$, проведите прямую.\\
«Функция возрастает» & Для $y=kx+b$ выполнено $k>0$.\\
«Пересекает $Oy$ ниже нуля» & Свободный коэффициент $b<0$.\\
«Точка пересечения графиков» & Приравняйте значения $y$, найдите $x$, затем $y$.\\
\bottomrule
\end{tabularx}
\end{center}
\clearpage

\begin{center}
{\Large\color{deepaccent}Тренировочный блок}\par
\vspace{2mm}
{\small Путь А --- 10 заданий, от простого к более сложному}
\end{center}
\vspace{2mm}
Выполняйте вычисления аккуратно. При построении подписывайте оси и точки. Для пересечения графиков проверяйте результат по обеим формулам.
\begin{enumerate}[label=\textbf{\arabic*.}]
\item Для $y=3x-2$ найдите $y(-1)$ и $y(2)$.
\item Для $y=-2x+3$ вычислите $y(2{,}5)$.
\item Для $y=-4x+1$ найдите $x$, если $y=-11$.
\item Для $y=1{,}5x-3$ найдите $x$, если $y=0$.
\item Постройте график $y=-x+2$ по таблице для $x=0$ и $x=2$. Укажите координаты обеих точек.
\item Для функции $y=\frac32x-2$ определите $k$, $b$, направление изменения функции и точку пересечения с $Oy$.
\item Сделайте эскиз прямой, у которой $k<0$ и $b>0$. Запишите один пример подходящей формулы.
\item Найдите точку пересечения графиков $y=2x+3$ и $y=-x+9$. Проверьте ответ по двум формулам.
\item Найдите точку пересечения графиков $y=4x-5$ и $y=-x+8$. Запишите координаты десятичными дробями и выполните проверку.
\item Даны три прямые: $y=2x-1$, $y=-x+5$ и $y=x+1$. Найдите точку пересечения первых двух и проверьте, принадлежит ли она третьей прямой.
\end{enumerate}
\clearpage

\begin{center}
{\Large\color{deepaccent}Ответы}\par
\vspace{2mm}
{\small Самопроверка после самостоятельного решения}
\end{center}
\vspace{6mm}
\begin{longtable}{@{}p{12mm}p{.78\linewidth}@{}}
\toprule
№ & Ответ\\
\midrule
\endhead
1 & $y(-1)=-5$, $y(2)=4$.\\
2 & $y(2{,}5)=-2$.\\
3 & $x=3$.\\
4 & $x=2$.\\
5 & $(0;2)$ и $(2;0)$; прямая убывает.\\
6 & $k=\frac32$, $b=-2$; возрастает; пересечение с $Oy$: $(0;-2)$.\\
7 & Убывает, пересекает $Oy$ выше нуля; например, $y=-x+2$.\\
8 & $(2;7)$; при $x=2$ обе формулы дают $y=7$.\\
9 & $(2{,}6;5{,}4)$; обе формулы дают $y=5{,}4$.\\
10 & $(2;3)$; третьей прямой принадлежит, поскольку $2+1=3$.\\
\bottomrule
\end{longtable}
\vfill
\begin{center}
\textcolor{softgray}{\small Лёвин Артём Александрович\quad$\cdot$\quad эксклюзивно для Кирилла Зиновьева}
\end{center}
\end{document}
